% Probeblatt der Vorlage mathblatt.sty (seit Stufe 6, Version 2026-09-28a)
% Zeigt jeden dokumentierten Baustein der Anleitung genau einmal, in der Reihenfolge der Anleitung,
% je Baustein mit seinem Namen in einer kleinen grauen Zeile darüber (\pbz). Inhalte aus den
% Beispielen der Anleitung, Grafiken verkleinert.
% Zweck: Lesestück für den Lehrer und Kompilierprobe jeder künftigen Stufe. Jeder Vorlagen-Auftrag
% ergänzt hier seine neuen Bausteine (an ihrer Stelle in der Anleitung) und prüft mit
%   python referenz/probeblatt-pruef.py
% dass Anleitung und Namenszeilen übereinstimmen. Zählgrenze: höchstens 12 Seiten.
% Kompilieren: mathblatt.sty neben diese Datei kopieren, zweimal xelatex (Abhakseite und
% Sprungziele). Das PDF liegt daneben im Repo (einzige Ausnahme von „Kompilate im Scratchpad").
% Seitenbausteine, die eine neue Seite beginnen (abhakseite, \abhakauto, \begleitteil,
% \hilfeseite), stehen am Ende, ihre Namenszeile am Fuß der Seite davor („nächste Seite");
% \blattkopf wirkt ab Seite 2, Seite 1 trägt die Fußzeile von \blattfuss. Mehrere Grafiken
% stehen nebeneinander (\pbfig); wo die Zeilenbreite es verlangt, weicht die Reihenfolge um
% wenige Plätze von der Anleitung ab (probeblatt-pruef.py nennt sie).
\documentclass[11pt]{article}
\usepackage{mathblatt}

% Namenszeile: \pbz{name1,name2,=umgebung} – eine kleine graue Zeile, je Name eine Randnotiz.
\newcounter{pbzahl}
\def\pbform#1#2\relax{\ifx=#1\textbackslash begin\{#2\}\else\textbackslash #1#2\fi}
\newcommand{\pbname}[1]{\stepcounter{pbzahl}\mbox{\pbform#1\relax}\hspace{1.2em}}
% (der Absatzabstand nach der Zeile fällt großteils weg, damit sie an ihrem Baustein hängt)
\newcommand{\pbz}[1]{\par\noindent{\scriptsize\color{mbgrau}\raggedright\foreach \pbn in {#1}{\expandafter\pbname\expandafter{\pbn}}\par}\nobreak\vskip-0.7\parskip}
% dasselbe für Bausteine, die selbst eine neue Seite beginnen: Zeile am Ende der Seite davor
\newcommand{\pbzn}[1]{\par\noindent{\scriptsize\color{mbgrau}\raggedright\foreach \pbn in {#1}{\expandafter\pbname\expandafter{\pbn}}(nächste Seite)\par}}
% Grafik mit Namenszeile in einer Spalte der Breite #1 (Anteil der Zeile)
\newcommand{\pbfig}[3]{\begin{minipage}[t]{#1\linewidth}\vspace{0pt}\pbz{#2}#3\end{minipage}}
% Abschnitt der Anleitung als Einheit
\newcommand{\pbteil}[3]{\einheitenkopf[#1]{Einheit #2 von 14 · #3}}

\begin{document}
\blattfuss{Probeblatt}{Vorlage Stufe 6}
\pbz{blattfuss,blattkopf}
{\small Seite 1 trägt die Fußzeile von \verb|\blattfuss| (unten links), ab Seite 2 gilt
\verb|\blattkopf*| mit Kopfzeile oben links (samt Einheit, Stufe 6) und Legende unten links.
Jede graue Zeile nennt die Bausteine, die direkt darunter stehen; Reihenfolge wie in der Anleitung.}

\pbz{einheitenkopf,zweigzeile}
\einheitenkopf*[e1]{Einheit 1 von 14 · Grundgerüst}
\zweigzeile{Den Rahmen eines Blatts kennen · Vorlage Stufe 6 · Kompilierprobe}
\pbz{verz,verzeichniszeile,verztrenn}
\verzeichniszeile{\verz{e1}{Grundgerüst} \verztrenn \verz{e2}{Wertetabelle} \verztrenn
  \verz{e3}{Koordinatensystem} \verztrenn \verz{e4}{3D} \verztrenn \verz{e5}{Körper} \verztrenn
  \verz{e6}{Kreis} \verztrenn \verz{e7}{Vierecke} \verztrenn \verz{e8}{Stochastik} \verztrenn
  \verz{e9}{Boxplot} \verztrenn \verz{e10}{Analysis} \verztrenn \verz{e11}{Zahlen} \verztrenn
  \verz{e12}{Trigonometrie} \verztrenn \verz{e13}{Prozent} \verztrenn \verz{e14}{schwach}
  \verztrenn \verz{abhaken}{Das kann ich}}

\pbz{uebersichtskasten}
\uebersichtskasten{Plusklammer: \quad $a + (b - c) = a + b - c$ \\ Minusklammer: \quad $a - (b - c) = a - b + c$}

\pbz{verfahren}
\verfahren{Teilaufgaben in einer Spalte und in zwei Spalten}
\begin{aufgabe}{Ich kann Teilaufgaben lesen. (Die Nummer bleibt mit allem darunter auf einer Seite.)}
\pbz{=aufgabe,=teile,teil,steil,=teilezwei,tz,leerfeld,feld,stz}
\begin{teile}
\teil $3 \cdot 4 = \leerfeld$
\steil $12 : 0{,}5 = \leerfeld$
\end{teile}
\begin{teilezwei}
\tz $y = 2x + 3$: \feld{m} & \stz $y = -x$: \feld{m} \\
\end{teilezwei}
\end{aufgabe}

\begin{aufgabe}{Ich kann Steigung und y-Achsenabschnitt ablesen.}
\pbz{=geruest,gz,gzs,=gleichungsraster,gl,sgl,anweisung}
\begin{geruest}
\gz{a}{$y = 3x + 2$}{\feld{m}\feld{n}}
\gzs{b}{$y = -x + 5$}{\feld{m}\feld{n}}
\end{geruest}
\begin{gleichungsraster}[2]
\gl{x + 5 = 9} & \sgl{7 - 2x = 15} \anweisung{Löse jetzt ohne Taschenrechner und mache die Probe.}
\gl[1]{2x + 3 = 11} & \\
\end{gleichungsraster}
\end{aufgabe}

\clearpage
\blattkopf*{Probeblatt}{Vorlage Stufe 6}{$\star$ = anspruchsvoll, darf übersprungen werden}

\begin{aufgabe}{Ich kann eine Gleichung senkrecht umformen.}
\pbz{rechenplatz,beispiel,rechnung}
\rechenplatz[halb]{2}
\beispiel{x + 5 &= 9 &&\mid -5 \\ x &= 4}
\rechnung{5x - 8 &= 12 &&\mid +8 \\ 5x &= 20}
\end{aufgabe}

\begin{aufgabe}{Ich kann am Streifen ablesen, wie viel Prozent ein Teil ist.}
\pbz{=beispiel,streifen}
\begin{beispiel}
Der Streifen sind 10 Kinder, gefärbt sind 4 Kinder. Wie viel Prozent sind gefärbt?
\streifen{40}{$0$}{10 Kinder}
\rechnung{1 \text{ Kästchen} &= 10\,\% \\ 4 \text{ Kästchen} &= 40\,\%}
Antwort: 40\,\% der Kinder sind gefärbt.
\end{beispiel}
\end{aufgabe}

\weit
\begin{aufgabe}{Ich kann Felder, Kästchen und Listen ausfüllen. (Diese Nummer steht unter \texttt{\textbackslash weit}.)}
\pbz{weit,feldl,punktfeld,janein,kreuz,mnliste,nullstellenliste,punktprobenliste}
\begin{teile}
\teil $y = 2x - 6$: \feldl{x_0} \quad Scheitelpunkt \punktfeld \quad Steigt die Gerade? \janein \quad \kreuz{A}\kreuz{B}
\end{teile}
\mnliste[2]{2x+3, -x}
\nullstellenliste{x-4}
\punktprobenliste{2x-1/A(3|5)}
\end{aufgabe}
\eng

\pbteil{e2}{2}{Wertetabelle}
\pbfig{0.55}{wertetabelle}{\wertetabelle[,4,,-2,-5]{x}{f(x)}{-2,-1,0,1,2}}\hfill
\pbfig{0.4}{wertetabelleleer}{\wertetabelleleer{x}{y}{4}}

\pbteil{e3}{3}{Koordinatensystem 2D}
\pbz{=ksys,gerade,punkt,steigungsdreieck,parabel,funktion,funktionab,ksysabstand}
\begin{ksys}[xmin=-3,xmax=4,ymin=-3,ymax=4,ablesen]
  \gerade{2}{-1}{g} \punkt{2}{3}{A} \steigungsdreieck{0}{-1}{2}
\end{ksys}\ksysabstand
\begin{ksys}[xmin=-3,xmax=4,ymin=-3,ymax=4,ablesen]
  \parabel{1}{-1}{-2}{p} \funktion{0.5*\x^2-2}{f} \funktionab{1/\x}{h}{0.25}{4}
\end{ksys}

\pbteil{e4}{4}{Koordinatensystem 3D}
\pbz{=ksys3,rpunkt,rvektor,rvektorab,rgerade,rebene,rebenepar,rquader,rpyramide}
\begin{ksys3}[x1min=0,x1max=2,x2min=0,x2max=4,x3min=0,x3max=3]
  \rebene{2}{4}{2}{E}
  \rpunkt{1}{2}{1}{A}
  \rgerade[-0.5:1]{0,1,2}{1,1,0}{g}
\end{ksys3}\hfill
\begin{ksys3}[x1min=0,x1max=2,x2min=0,x2max=4,x3min=0,x3max=3]
  \rvektor{0,3,2.5}{\vec a}
  \rvektorab[below]{2,0,0}{2,2,0}{\vec u}
  \rebenepar*{0,0,2}{1,0,0}{0,1,0}{D}
  \rquader{0,3,0}{1}{1}{1}
  \rpyramide{1,1.5,0}{1}{1}{1.5,2,2}
\end{ksys3}
\pbteil{e5}{5}{Geometrie und Körper}
\pbfig{0.24}{dreieckrw}{\dreieckrw{3}{2}{a}{b}{c}}\hfill
\pbfig{0.27}{dreieck}{\dreieck{(0,0)}{(3.5,0)}{(1,2)}{a}{b}{c}{\alpha}{\beta}{\gamma}}\hfill
\pbfig{0.22}{quader}{\quader{2.5}{1.5}{1.8}{a}{b}{c}}\hfill
\pbfig{0.18}{zylinder}{\zylinder{0.8}{2}{r}{h}}

\medskip
\pbfig{0.27}{prismadreieck}{\prismadreieck{2.5}{1.8}{1.5}{g}{h}{l}}\hfill
\pbfig{0.25}{pyramide}{\pyramide{2.5}{2}{2.5}{a}{b}{h}}\hfill
\pbfig{0.2}{kegel}{\kegel{1}{2}{r}{h}{s}}\hfill
\pbfig{0.18}{kugel}{\kugel{1}{r}}

\pbteil{e6}{6}{Kreis und Winkelfiguren}
\pbfig{0.3}{=kreis,mittelpunkt,kreispunkt,radius,durchmesser,sehne}{%
\begin{kreis}[1.2]
  \mittelpunkt{M} \kreispunkt{150}{B} \radius{70}{r} \durchmesser{0}{d} \sehne{200}{300}{s}
\end{kreis}}\hfill
\pbfig{0.3}{tangente,sektor,bogen}{%
\begin{kreis}[1.2]
  \mittelpunkt{} \tangente{300}{t} \sektor{30}{110}{\alpha} \bogen{150}{230}{b}
\end{kreis}}\hfill
\pbfig{0.36}{mittelpunktswinkel,umfangswinkel}{%
\begin{kreis}[1.2]
  \mittelpunkt{} \mittelpunktswinkel{200}{340}{\mu} \umfangswinkel{200}{340}{100}{\gamma}
\end{kreis}}

\medskip
\pbfig{0.29}{geradenkreuzung}{\geradenkreuzung{35}{\alpha}{\beta}{\gamma}{\delta}}\hfill
\pbfig{0.26}{parallelenpaar}{\parallelenpaar{60}{\alpha}{\beta}{\gamma}{\delta}}\hfill
\pbfig{0.13}{winkel}{\winkel[2]{40}{\alpha}}\hfill
\pbfig{0.28}{winkelstrahl}{\winkelstrahl[3.5]}

\pbteil{e7}{7}{Vierecke, Netze, Strahlensatz}
\pbfig{0.32}{viereck}{\viereck[seiten={a,b,c,d},diagonalen]{(0,0)}{(3.5,0)}{(4,2)}{(1,2.4)}}\hfill
\pbfig{0.32}{parallelogramm}{\parallelogramm[seiten={a,b,a,b},hoehe]{3}{2}{60}}\hfill
\pbfig{0.3}{rechteck}{\rechteck[achsen=beide]{3}{1.8}}

\medskip
\pbfig{0.32}{trapez}{\trapez[hoehe=h]{3.5}{2}{1.8}}\hfill
\pbfig{0.32}{raute}{\raute[achsen=diagonalen,diagonalen]{3}{2}}\hfill
\pbfig{0.3}{drachen}{\drachen[achsen=senkrecht,diagonalen]{3}{2.4}{0.35}}

\medskip
\pbfig{0.4}{netzquader}{\netzquader{1.5}{1}{1}{a}{b}{c}}\hfill
\pbfig{0.25}{netzwuerfel}{\netzwuerfel{0.9}{a}}\hfill
\pbfig{0.3}{netzpyramide}{\netzpyramide{1.5}{1.4}{a}{h_s}}

\medskip
\pbfig{0.4}{netzzylinder}{\netzzylinder{0.5}{1.5}{r}{h}}\hfill
\pbfig{0.55}{strahlensatz}{\strahlensatz[strecken={2,3,,5}]{1.6}{4}{35}}

\pbteil{e8}{8}{Stochastik}
\pbfig{0.48}{baumzwei}{\baumzwei{R/0{,}4,B/0{,}6}{R/0{,}4,B/0{,}6}{R/,B/}}\hfill
\pbfig{0.48}{saeulen}{\saeulen{Mo/12,Di/7,Mi/9}{16}{4}{Anzahl}}

\medskip
\pbfig{0.48}{saeulenab}{\saeulenab[ymax=100,ystep=20,yfein=5,ylabel=Gäste,hoehe=3.5]{Mo/45,Di/70,Mi/30}}\hfill
\pbfig{0.48}{balkenab}{\balkenab[xmax=40,xstep=10,xfein=2,xlabel=Nennungen,breite=5,hoehe=0.8]{Rot/24,Blau/32,Grün/14}}

\medskip
\pbfig{0.48}{liniendia}{\liniendia[ymax=20,ystep=5,yfein=1,ylabel=Temperatur in $^\circ$C,hoehe=3]{Jan/2,Feb/5,Mär/4,Apr/11}}\hfill
\pbfig{0.48}{baumdreigleich}{\baumdreigleich{R/0{,}4,B/0{,}6}}

\medskip
\pbfig{0.45}{baumdrei}{\baumdrei{R/\frac{2}{5},B/\frac{3}{5}}{R/\frac{1}{4},B/\frac{3}{4}}{R/,B/}{R/,B/}{R/,B/}{R/,B/}{R/,B/}}\hfill
\pbfig{0.27}{kreisdiagramm}{\kreisdiagramm[0.9]{Bus/40, Rad/25, Auto/20, Fuß/15}}\hfill
\pbfig{0.11}{kreisleer}{\kreisleer[0.9]}\hfill
\pbfig{0.12}{kreissektor}{\kreissektor[0.9]{135}{135°}}

\medskip
\pbfig{0.35}{kreisdiagrammleer}{\kreisdiagrammleer[0.9]{Bus/40, Rad/25, Auto/35}}\hfill
\pbfig{0.4}{sachtabelle,leerzelle,strichliste}{\sachtabelle{lcc}{Medium & M & J}{Handy & 97\,\% & \leerzelle\\ Striche & \strichliste{13} & 32}}

\medskip
\pbfig{0.39}{vierfeldertafel}{\vierfeldertafel{A}{B}{20,30,50,10,40,50,30,70,100}}\hfill
\pbfig{0.32}{binomialverteilung}{\binomialverteilung[von=0,bis=6]{10}{0.3}}\hfill
\pbfig{0.24}{normalverteilung}{\normalverteilung[von=-1,bis=1,karo=0.45]{0}{1}}
\pbteil{e9}{9}{Boxplot und Histogramm}
\pbfig{0.5}{=boxplots,bp}{%
\begin{boxplots}[xmin=0,xmax=20,xstep=2,xlabel=Punkte]
  \bp{4,7,9,13,18}{Klasse 8a}
  \bp{}{Klasse 8b}
\end{boxplots}}\hfill
\pbfig{0.46}{histogramm}{\histogramm[ymax=10,ystep=2,xstep=10,ylabel=$h$]{0:10/4, 10:20/9, 20:30/6}}

\pbteil{e10}{10}{Analysis}
\pbfig{0.4}{ableitungspaar}{\ableitungspaar[xmin=-2,xmax=2,ymin=-1,ymax=3,ymin2=-2,ymax2=2,klein]{0.5*\x^2-1}{f}{\x}{f'}}\hfill
\pbfig{0.56}{flaeche,flaechezwischen,tangentean,hochpunkt,tiefpunkt,wendepunkt,asymptote}{%
\begin{ksys}[xmin=-2,xmax=4,ymin=-2,ymax=4,ablesen]
  \funktion{\x^2-4*\x+3}{}
  \flaeche{\x^2-4*\x+3}{1}{3}{A}
  \flaechezwischen{-\x^2+2*\x+3}{\x+1}{-1}{2}{B}
  \tangentean{0.5*\x^2}{1}{t}
  \hochpunkt{1}{4}{H} \tiefpunkt{2}{-1}{T} \wendepunkt{0}{0}{W}
  \asymptote{x=3.5}{}
\end{ksys}}

\pbteil{e11}{11}{Zahlen und Algebra}
\pbz{zahlenstrahl}
\zahlenstrahl[xmin=7.6,xmax=8,xstep=0.01,karo=0.4]{7.68/7{,}68}

\pbfig{0.39}{=zahlengerade,intervall,intervallo}{%
\begin{zahlengerade}[xmin=-4,xmax=6,karo=0.6]
  \intervall{-2}{3}{[-2;3]}
  \intervallo[2]{1}{5}{]1;5[}
  \intervallo*[3]{2}{}{x \ge 2}
\end{zahlengerade}}\hfill
\pbfig{0.12}{bruchkreis}{\bruchkreis[0.9]{3}{8}}\hfill
\pbfig{0.18}{bruchrechteck}{\bruchrechteck[3]{3}{8}}\hfill
\pbfig{0.26}{termbaum,tb}{\termbaum{+}{3x}{\tb{\cdot}{2}{y}}}

\pbteil{e12}{12}{Trigonometrie}
\pbfig{0.3}{einheitskreis}{\einheitskreis[karo=1.3]{50}}\hfill
\pbfig{0.66}{sinus,kosinus}{%
\begin{ksys}[trigo,karo=0.7]
  \sinus{1}{1}{f} \kosinus{2}{0.5}{g}
\end{ksys}}

\pbteil{e13}{13}{Prozentrechnung: Streifen und Dreisatz}
\pbz{streifenleer}
\streifenleer[0]

\pbz{streifenvoll}
\streifenvoll{Rad/40,Bus/25,zu Fuß/20,Auto/15}

\pbz{streifenfrage}
\streifenfrage{65}{$0\,\%$}{$100\,\%$}
\pbz{streifenfeld}
\streifenfeld[4]{75}{$0\,\%$}{$100\,\%$}
\begin{minipage}[t]{0.66\linewidth}\vspace{0pt}
\pbz{streifenreihe}
\streifenreihe{a/20,b/60}
\pbz{streifenwertreihe}
\streifenwertreihe{a/30/6\,kg/, b/70//50\,€}
\end{minipage}\hfill
\begin{minipage}[t]{0.3\linewidth}\vspace{0pt}\pbz{=dreisatz,dsz,dsp,dsleer}
\begin{dreisatz}{Hefte}{Preis}
\dsz{4}{6{,}00\,€}
\dsp{:4}{:4}
\dsz{1}{\dsleer}
\dsp{\cdot 6}{\cdot 6}
\dsz{6}{9{,}00\,€}
\end{dreisatz}
\end{minipage}

\pbteil{e14}{14}{Option schwach}
\begin{aufgabe}{Ich kann den Anteil am Streifen ausrechnen. (Option schwach)}
\pbz{swz,swa,swb,swfrage}
\swb{\beispiel{20 \text{ von } 100 &= \tfrac{20}{100} = 20\,\%}}{\bruchrechteck[5]{1}{5}}
\swz[1]{Der Streifen sind 20\,€, gefärbt sind 6\,€.}{\streifen{30}{$0\,€$}{$20\,€$}}
\swa{Ist mehr als die Hälfte gefärbt?}{\kreuz{ja}\kreuz{nein}}
\swfrage{Was bleibt gleich, was ändert sich?}
\end{aufgabe}

\pbzn{=abhakseite,abhakgruppe,abhak}
\begin{abhakseite}[Das kann ich (Handform)]
\zweigzeile{Hake ab, was du kannst.}
\abhakgruppe{Einheit 1 · Grundgerüst}
\abhak{1}{Ich kann Teilaufgaben lesen.}
\abhak{Z1}{Ich kann einen langen Titel lesen, der umbricht und unter dem Titel einhängt, nicht unter dem Kästchen, auch wenn die Nummer breiter ist.}
\end{abhakseite}

\pbzn{abhakauto}
\abhakauto

\pbzn{begleitteil,erg}
\begleitteil
\erg{1}{a) 12 \quad b) 24 \quad c) $m = 2$ \quad d) $m = -1$}

\pbzn{hilfeseite,=schritte,schritt,achtung}
\hilfeseite
\verfahren{Gleichung lösen}
\begin{schritte}
\schritt Klammern auflösen.
\schritt Alles mit $x$ auf eine Seite bringen.
\achtung{Beim Teilen durch eine negative Zahl das Vorzeichen beachten.}
\end{schritte}

\bigskip\noindent{\small\color{mbgrau}Randnotizen auf diesem Probeblatt: \thepbzahl}
\end{document}
